\documentclass[11pt]{article}
\usepackage[margin=1in]{geometry}
\usepackage{xcolor}
\usepackage{enumitem}
\usepackage{amsmath}
\usepackage{amssymb}
\usepackage{hyperref}
\hypersetup{colorlinks=true,linkcolor=black,urlcolor=blue}
\definecolor{revblue}{RGB}{0,0,180}
\newcommand{\loc}[1]{\textbf{[\textcolor{revblue}{Tracked PDF: #1}]}}
\setlist{nosep,leftmargin=1.4em}
\setlength{\parskip}{0.55em}
\setlength{\parindent}{0pt}

\begin{document}

\begin{center}
{\Large\textbf{Cover Letter to Reviewer 2}}\\[0.5em]
\textbf{Round 2 Revision}
\end{center}

\textbf{Manuscript:} ``Content-Conditioned Invariant Feature-Statistic
Augmentation: Unifying Robustness and Weight Averaging for Domain
Generalization'' (CIFA)\\
\textbf{Journal:} \textit{Electronics} (MDPI)\qquad
\textbf{Manuscript ID:} electronics-4654208\\
\textbf{Date:} 3 November 2026

\hrule
\vspace{0.4em}

Dear Reviewer 2,

We thank you for the candid and careful assessment, and for acknowledging that
the presentation issues have been addressed, that the claims have been
appropriately revised, and that the null and non-replicated results are reported
transparently. The remaining concerns concern the interpretation and
positioning of the work rather than additional experiments. We respond to each
point below; the location tag at the start of every response gives the exact
page and line range in the tracked PDF, which uses continuous line numbering.

\subsection*{Point 1 --- Necessity of the distinctive components}
\loc{Page 33, Lines 1070--1074; Proposition 1, Pages 8--9, Lines 304--348}

We agree that the content--style decomposition is not a source of measured
accuracy: full-moment mixing is statistically indistinguishable from full CIFA
(paired $+0.23\pm0.75$, $p=0.53$), and removing weight averaging is the
dominant ablation ($2.42$ points). We would, however, distinguish accuracy
contribution from theoretical role. The decomposition is not intended to
improve accuracy; it is the structural condition that makes the
content-preservation bound admissible. By mixing only the style deviation and
fixing the content baseline, the construction guarantees
$\mathbb{E}[\mathbf{d}\mid Y]=0$ and thereby licenses the total-variation bound
of Proposition~1. Mixing the full moments does not satisfy this condition and
cannot carry the same guarantee, even when its accuracy is matched. The two
constructions are tied in accuracy but not in the guarantees they provide;
empirical equivalence in accuracy does not establish logical equivalence in the
conditions each satisfies.

\subsection*{Point 2 --- Value of the unification}
\loc{Pages 8--11, Lines 304--400}

We explicitly position the contribution as integration rather than a new
operator. The unification is not a concatenation of existing techniques: it
places feature-moment augmentation, convex-hull minimax, and weight averaging
under one causal-invariance principle and states the conditions under which
each applies. This yields testable predictions---content--style separability as
an enabling rather than sufficient condition, the behavior of the coverage
gap, and the second-order flat-minima approximation. Together with the
reproducible platform and the falsifiable boundary, this allows later methods
to be compared on a common footing and avoids repeated investment in directions
already shown to be ineffective. We do not claim a deployment advantage; the
practical significance lies in the unified understanding and the
well-characterized baseline.

\subsection*{Point 3 --- Strength of the assumptions and the coverage term}
\loc{Page 8, Lines 304--322; Page 36, Lines 1171--1173}

The assumptions are now stated, tested, and falsifiable premises rather than
hidden conditions. Content--style separability is labeled an enabling premise,
and the trajectory--augmentation alignment is presented as an empirically
supported assumption rather than a theorem, with its limitations stated. The
target access required to evaluate the coverage-gap term is an inherent
property of target-free generalization, not a limitation unique to this work; we
demonstrate the behavior of both terms in a controlled moment-space setting
(Figure~2, Page~14) rather than presenting the term as computable without
target data. A strong assumption is not an incorrect one; the value of the
analysis is that its conditions are explicit and open to falsification.

\subsection*{Point 4 --- Weakening claims versus adding evidence}
\loc{Pages 21--23, Lines 795--828; Pages 8--11, Lines 304--400}

When the evidence does not support a stronger claim, weakening that claim is
the scientifically correct response; retaining an unsupported claim would be
the error. The revision also adds positive evidence---an independent
unified-code replication, a gate-pressure stress test, the alignment
verification, and a corrected two-one-sided equivalence analysis. This
evidence points to statistical equivalence and to boundary conditions rather
than to superiority. What was weakened were the claims that did not hold; the
claims that do hold---the unifying framework and the boundary
characterization---are retained, strengthened, and made more precise and more
readily falsifiable.

\subsection*{Point 5 --- Overall scientific significance}
\loc{Page 36, Lines 1168--1200}

The manuscript does not claim a performance breakthrough. Its scientific
significance is at the level of understanding: a unifying explanation, a
systematic validity boundary, and a falsifiable, reproducible benchmark; the
transparent reporting of null and non-replicated results is itself part of that
contribution. The only actionable editorial item has been completed (Point~6).
We respectfully ask that the work be judged on what it actually establishes
rather than on the overstated claims of the original submission.

\subsection*{Point 6 --- Editorial cleanup (duplicated or partially edited text)}
\loc{Page 37, Lines 1209--1211}

We carried out a full editorial review. The one concrete duplication was the
data-availability statement, which appeared twice: a manually typeset
``\textbf{Data availability.}'' paragraph and the MDPI template
\texttt{\textbackslash dataavailability\{\}} command with identical content.
The manual paragraph has been removed and the template command retained, so the
statement now renders exactly once. The remaining items that can look like
duplication in a text extraction (redefined revision macros, deliberate
indentation and table alignment) were inspected individually and are not
duplicated content.

\vspace{0.6em}
We thank you once again for the thorough and transparent review.

\vspace{0.8em}
Sincerely,\\
Qingsong Cai (on behalf of all authors)\\
\textit{Beijing Technology and Business University}\qquad
\textit{caiqs@btbu.edu.cn}

\end{document}
